% GENERATED FILE. Edit canonical lessons, then run npm run generate:books.
% canonical-source-hash: fnv1a64:b33a2d486643c063
% canonical-lessons: LA-C127-faba, LA-C127-pisum, LA-C127-cucumis, LA-C127-pullus, LA-C127-piper

\chapter{Bean, Pea, Cucumber, Chicken, Pepper}
\label{ch:la-bean-pea-cucumber-chicken-pepper}
\hlchaptermodality{127}

\begin{chapteropening}
\textbf{By the end of this chapter:} \emph{I can say a bean, a pea, a cucumber, a chicken and pepper.}

Complete the last lesson of chapter 127: I can say a bean, a pea, a cucumber, a chicken and pepper.
\end{chapteropening}

\section[faba, fabae]{faba, fabae — a bean}

\label{lesson:LA-C127-faba}



\begin{quote}
Before the new one: say the Latin for an onion, then the Latin for garlic.
\end{quote}

\subsection*{You'll want to know: faba, fabae}
\textbf{faba, fabae} — \textquotedblleft{}a bean\textquotedblright{}.

\begin{cousinweb}
\textbf{faba} is inherited, usually traced to Proto-Indo-European \emph{*bhabh-}, \textquotedblleft{}bean\textquotedblright{} (EtymDB 2.1). English \emph{bean} may be a distant cousin.
\end{cousinweb}

\begin{grammarlens}[title={What you've built}]
One more word for this chapter.
\end{grammarlens}

\subsection*{Your turn}
\begin{itemize}
\raggedright
  \item \emph{Say it:} \emph{faba}
  \item \emph{Say it:} \emph{faba}, once more
  \item \emph{Say it:} say \emph{faba}
  \item \emph{Recall:} say the Latin for an onion, then the Latin for garlic, then say \emph{faba} again
\end{itemize}

\begin{tcolorbox}[breakable,colback=teal!4,colframe=teal!35!black,title={Before you move on}]
What does faba, fabae mean? (\textquotedblleft{}A bean\textquotedblright{}.) Say it once more.
\end{tcolorbox}
\section[pīsum, pīsī]{pīsum, pīsī — a pea}

\label{lesson:LA-C127-pisum}



\begin{quote}
Before the new one: say the Latin for garlic, then the Latin for a bean.
\end{quote}

\subsection*{You'll want to know: pīsum, pīsī}
\textbf{pīsum, pīsī} — \textquotedblleft{}a pea\textquotedblright{}.

\begin{cousinweb}
\textbf{pīsum} was borrowed from Greek \emph{pison}, \textquotedblleft{}pea\textquotedblright{} (Lewis \& Short). English \emph{pea} comes from it: older English \emph{pease} was heard as a plural.
\end{cousinweb}

\begin{grammarlens}[title={What you've built}]
One more word for this chapter.
\end{grammarlens}

\subsection*{Your turn}
\begin{itemize}
\raggedright
  \item \emph{Say it:} \emph{pīsum}
  \item \emph{Say it:} \emph{pīsum}, once more
  \item \emph{Say it:} say \emph{pīsum}
  \item \emph{Recall:} say the Latin for garlic, then the Latin for a bean, then say \emph{pīsum} again
\end{itemize}

\begin{tcolorbox}[breakable,colback=teal!4,colframe=teal!35!black,title={Before you move on}]
What does pīsum, pīsī mean? (\textquotedblleft{}A pea\textquotedblright{}.) Say it once more.
\end{tcolorbox}
\section[cucumis, cucumeris]{cucumis, cucumeris — a cucumber}

\label{lesson:LA-C127-cucumis}



\begin{quote}
Before the new one: say the Latin for a bean, then the Latin for a pea.
\end{quote}

\subsection*{You'll want to know: cucumis, cucumeris}
\textbf{cucumis, cucumeris} — \textquotedblleft{}a cucumber\textquotedblright{}.

\begin{cousinweb}
\textbf{cucumis} has no sure origin; it is probably a loan from a lost Mediterranean language. English \emph{cucumber} comes from it through French.
\end{cousinweb}

\begin{grammarlens}[title={What you've built}]
One more word for this chapter.
\end{grammarlens}

\subsection*{Your turn}
\begin{itemize}
\raggedright
  \item \emph{Say it:} \emph{cucumis}
  \item \emph{Say it:} \emph{cucumis}, once more
  \item \emph{Say it:} say \emph{cucumis}
  \item \emph{Recall:} say the Latin for a bean, then the Latin for a pea, then say \emph{cucumis} again
\end{itemize}

\begin{tcolorbox}[breakable,colback=teal!4,colframe=teal!35!black,title={Before you move on}]
What does cucumis, cucumeris mean? (\textquotedblleft{}A cucumber\textquotedblright{}.) Say it once more.
\end{tcolorbox}
\section[pullus, pullī]{pullus, pullī — a chicken, a young animal}

\label{lesson:LA-C127-pullus}



\begin{quote}
Before the new one: say the Latin for a pea, then the Latin for a cucumber.
\end{quote}

\subsection*{You'll want to know: pullus, pullī}
\textbf{pullus, pullī} — \textquotedblleft{}a chicken, a young animal\textquotedblright{}.

\begin{cousinweb}
\textbf{pullus} is inherited, usually traced to Proto-Indo-European \emph{*polH-}, \textquotedblleft{}young animal\textquotedblright{} (EtymDB 2.1); Greek \emph{pōlos}, \textquotedblleft{}foal\textquotedblright{}, and English \emph{foal} are cousins (Lewis \& Short). English \emph{pullet} comes from \emph{pullus} through French.
\end{cousinweb}

\begin{grammarlens}[title={What you've built}]
One more word for this chapter.
\end{grammarlens}

\subsection*{Your turn}
\begin{itemize}
\raggedright
  \item \emph{Say it:} \emph{pullus}
  \item \emph{Say it:} \emph{pullus}, once more
  \item \emph{Say it:} say \emph{pullus}
  \item \emph{Recall:} say the Latin for a pea, then the Latin for a cucumber, then say \emph{pullus} again
\end{itemize}

\begin{tcolorbox}[breakable,colback=teal!4,colframe=teal!35!black,title={Before you move on}]
What does pullus, pullī mean? (\textquotedblleft{}A chicken, a young animal\textquotedblright{}.) Say it once more.
\end{tcolorbox}
\section[piper, piperis]{piper, piperis — pepper}

\label{lesson:LA-C127-piper}



\begin{quote}
Before the new one: say the Latin for a cucumber, then the Latin for a chicken.
\end{quote}

\subsection*{You'll want to know: piper, piperis}
\textbf{piper, piperis} — \textquotedblleft{}pepper\textquotedblright{}.

\begin{cousinweb}
\textbf{piper} was borrowed from Greek \emph{peperi}, \textquotedblleft{}pepper\textquotedblright{} (Lewis \& Short), which came by trade from an Indian word like Sanskrit \emph{pippalī}, \textquotedblleft{}long pepper\textquotedblright{} (EtymDB 2.1).
\end{cousinweb}

\begin{grammarlens}[title={What you've built}]
One more word for this chapter.
\end{grammarlens}

\subsection*{Your turn}
\begin{itemize}
\raggedright
  \item \emph{Say it:} \emph{piper}
  \item \emph{Say it:} \emph{piper}, once more
  \item \emph{Say it:} say \emph{piper}
  \item \emph{Recall:} say the Latin for a cucumber, then the Latin for a chicken, then say \emph{piper} again
\end{itemize}

\begin{tcolorbox}[breakable,colback=teal!4,colframe=teal!35!black,title={Before you move on}]
What does piper, piperis mean? (\textquotedblleft{}Pepper\textquotedblright{}.) Say it once more.
\end{tcolorbox}
